\documentclass[11pt]{article}

\usepackage[margin=25mm]{geometry}

\usepackage{amsmath,amssymb,fontspec,unicode-math}

\IfFontExistsTF{DejaVu Serif}{\setmainfont{DejaVu Serif}}{\setmainfont{Latin Modern Roman}}

\IfFontExistsTF{Latin Modern Math}{\setmathfont{Latin Modern Math}}{\setmathfont{DejaVu Math TeX Gyre}}

\usepackage{xurl,hyperref}

\setlength{\parindent}{0pt}\setlength{\parskip}{6pt}

\begin{document}

\title{Role Bias in Diffusion Models: Diagnosing and Mitigating through Intermediate Decomposition}

\author{ELVIA preliminary AI analysis}\date{21 October 2026}\maketitle

Preliminary assessment: Minor.\par

Independent scientific review remains pending. Author adoption has not been requested.

Single-component flags in the later SDXL release mask logs but retain both optimized losses. Its role-difference score supports relative preference rather than absolute alignment, and a finite intermediate limit causes inconsistent dictionary access. Original ablation correspondence remains unresolved.

\section*{Review scope}

Claim-focused reading of the original paper and relevant author-released source, with direct mathematical and static-source checks. No author program, historical verifier, model training or published benchmark was executed.

\section*{Problem 1: Single-component flags do not mask the optimized supervised loss}

Source: Section 3.3, Eq.(3), PDF p.5; Table 3, p.9; rebind/scripts/sdxl/rebind\_lora\_sdxl.py, lines 1489–1507

Assessment: supported; impact: minor

\subsection*{Source claim (paraphrase)}

In the ordinary no-SNR branch, active/passive selection flags zero the corresponding logging variables but do not affect the weighted loss returned for optimization.

\subsection*{Analysis}

Use one active and one passive residual \(r_A=1\), \(r_P=3\), with weights \(5\) and \(1\). The returned mean is \((5+9)/2=7\) under active-only, passive-only and combined flags. The derivative with respect to the allegedly disabled passive residual is \(r_P=3\), whereas a genuine active-only loss has derivative zero. Both source groups are loaded; the combined method itself legitimately uses both.

\subsection*{Finding and scope}

The advertised public flag cannot by itself reproduce a single-component ablation. The code revision or external data selection used for Table 3 is not established.

\subsection*{Proposed correction}

Apply component masks to the optimized loss itself, for example \(A\lambda n_A^{-1}\sum r_{A,i}^2+P n_P^{-1}\sum r_{P,i}^2\), omitting disabled terms and their loader requirements. Log the same included terms.

A proposed correction is described; author adoption is not established.

\section*{Problem 2: Relative role preference is not an absolute alignment threshold}

Source: Section 4, Bias paragraph, PDF p.6; Table 3 and interpretation, p.9; load\_triplet\_data, lines 859–928

Assessment: supported; impact: minor

\subsection*{Source claim (paraphrase)}

The paper defines \(\beta=U-M\) and interprets \(\beta\le0\) as an unbiased, well-aligned image. The public pruning path tests relative ordering and does not use its absolute-threshold argument.

\subsection*{Analysis}

The criterion is exactly \(M\ge U\), satisfied by \((M,U)=(0.01,0)\) or \((0,0)\); it gives no positive lower bound on matching score. At combined weight \(1\), Table 3 gives matching/reversed scores \(63.4,75.5\); at weight \(5\), they are \(57.9,65.8\). The reported direction score improves by \(4.1\) percentage points while matching score falls by \(5.5\). Thus reduced reverse preference is supported, but the stated absolute-matching interpretation does not follow from that statistic. The public filter retains zero and low-score ties because only the sign of \(M-U\) is tested.

\subsection*{Finding and scope}

The relative statistic remains useful within its definition. No real image is assigned the synthetic scores, and human or fine-grained evaluation evidence is not rejected.

\subsection*{Proposed correction}

Describe \(\beta\le0\) as intended-over-reversed score ordering and report both scores. If absolute quality is intended, separately calibrate an explicit condition such as \(M\ge\tau\) together with \(M-U\ge\delta\); no numerical threshold is validated here.

A proposed correction is described; author adoption is not established.

\section*{Problem 3: A finite intermediate limit is applied inconsistently across two passes}

Source: create\_intermediate\_dpo\_dataset, rebind/scripts/sdxl/rebind\_lora\_sdxl.py, lines 1017–1027 and 1083–1094

Assessment: supported; impact: minor

\subsection*{Source claim (paraphrase)}

The first pass slices each relation-type list to the requested finite limit, but the construction pass iterates through the unsliced list and indexes the partially populated dictionary.

\subsection*{Analysis}

With active relations \((a_1,a_2)\), passive relations \((p_1,p_2)\) and limit \(1\), loading creates only the first relation of each type and its reverse. Construction nevertheless reaches \(a_2\) and accesses its absent dictionary entry. The full-selection default loads all those keys and is a passing control; an image-count ablation is a different parameter.

\subsection*{Finding and scope}

The nondefault finite-relation configuration has a direct missing-key path. This does not show that a default experiment crashed or that Table 10’s image-count experiment is invalid.

\subsection*{Proposed correction}

Compute the selected relation lists once and reuse them during loading and construction. This prevents requests for unselected relation keys, without guaranteeing that all unrelated dataset errors are resolved.

A proposed correction is described; author adoption is not established.

\section*{Limits}

Preliminary AI assessment; independent human scientific adjudication remains pending.\par

Analytic examples and release discrepancies do not remeasure the paper’s reported experiments.\par

The February 2026 public SDXL snapshot postdates the conference; correspondence to original generating runs is unestablished.\par

\section*{Source identity}

arXiv:2503.10037v3 (current author version confirmed from official metadata)

\url{https://arxiv.org/pdf/2503.10037v3}

Reviewed PDF SHA-256: \nolinkurl{b0a11d5a4d2137ba00ebe812236d0c37e831c7c4628a8e1b10366c04dd0ca67d}

The SHA-256 identifies the exact retained PDF inspected in this rapid source/argument check. Original supplied report bytes are retained locally. No unambiguous original PDF retrieval timestamp was recovered; file modification time was not substituted. Current official arXiv metadata checked at 2026-10-21T05:36:35.986302+00:00 identifies 2503.10037v3; implicated passages were checked in the current identified material.

\end{document}
